\section{Methodology}
\subsection{Preliminaries}
\label{sec:preliminaries}

\paragraph{Transformer Language Models.}
We consider a standard causal Transformer-based language model parameterized by $\theta$. Given an input token sequence
$\mathbf{x} = (x_1, \dots, x_T)$, each token is mapped to an embedding $e(x_t) \in \mathbb{R}^d$. Let
\[
E_t = [e(x_1), \dots, e(x_t)]
\]
denote the input embedding sequence up to position $t$. The Transformer computes a sequence of hidden states
\[
H_t = \mathrm{Transformer}_\theta(E_t) \in \mathbb{R}^{t \times d},
\]
where $h_t = H_t[t]$ is the final hidden state at position $t$. Next-token prediction is performed via a linear output head
\[
p(x_{t+1} \mid x_{\le t}) = \mathrm{softmax}(W h_t),
\]
with $W \in \mathbb{R}^{|\mathcal{V}| \times d}$ and vocabulary $\mathcal{V}$. The model is trained using standard
autoregressive maximum likelihood with a cross-entropy loss.

\paragraph{Chain-of-Thought and Latent Reasoning.}
In chain-of-thought (CoT) reasoning, the model explicitly generates intermediate reasoning steps as language tokens before producing the final answer. While effective, this constrains reasoning to the discrete language space and requires autoregressive generation of each reasoning token, tying both computation and representational structure to the surface form of language.

\paragraph{Chain of Continuous Thought (Coconut).}
Coconut proposes to shift reasoning from the language space to a continuous latent space. Instead of decoding intermediate hidden states into language tokens, Coconut directly feeds hidden states back into the model as subsequent input embeddings, enabling the model to ``think'' in a continuous space. Conceptually, this mechanism resembles a recurrent state update: the hidden representation at one reasoning step is reused as the input state for the next step, allowing information to be accumulated and refined across iterations. Unlike classical recurrent neural networks (RNNs), however, Coconut implements this recurrence using a Transformer backbone and leverages pretrained language representations.

Formally, Coconut introduces a latent reasoning mode delimited by special tokens \texttt{<start-thinking>} and \texttt{<end-thinking>}. When the model is in latent mode, the input embedding at position $t$ is defined as
\[
E_t = [e(x_1), \dots, e(x_i), h_i, h_{i+1}, \dots, h_{t-1}],
\]
where $x_i = \texttt{<start-thinking>}$ and $h_{t-1}$ is the hidden state from the previous position. In this regime, the hidden state $h_{t-1}$ serves as a \emph{continuous thought} that represents the current reasoning state. No language token is generated for latent positions, and the language modeling loss is masked over these steps.

\paragraph{Training Procedure in Coconut.}
make it small
Coconut is trained using a multi-stage curriculum that leverages supervised CoT data. At the initial stage, the model is trained on full language-based reasoning chains. In subsequent stages, the first $k$ language reasoning steps are progressively replaced by $k \times c$ continuous thoughts, where $c$ is a hyperparameter controlling the number of latent steps per removed language step. The training objective remains standard next-token prediction on the remaining language tokens, with losses masked over latent positions. This curriculum encourages the model to internalize reasoning processes into continuous latent representations while remaining compatible with pretrained language models.

---

\subsection{Limitations of Sequential Continuous Thought}
\label{sec:limitations}

Despite its conceptual appeal, Coconut introduces several structural limitations stemming from its sequential latent generation mechanism.

\paragraph{Sequential Latent Unrolling.}
In Coconut, each continuous thought is generated by feeding the hidden state from the previous step back into the model as the next input embedding. Consequently, generating $n$ latent reasoning steps requires $n$ sequential forward passes through the Transformer, followed by an additional pass to compute the loss on language tokens. This design tightly couples the depth of latent reasoning to sequential execution, resulting in inherently limited parallelism.

This sequential unrolling effectively reintroduces a core limitation that motivated the transition from recurrent architectures to Transformers. One of the primary advantages of Transformers over RNNs is their ability to process sequence elements in parallel, avoiding the step-by-step recurrence that hinders scalability and efficient hardware utilization. By enforcing a recurrence over latent reasoning steps, Coconut forfeits this parallelism at the reasoning level: each latent update depends on the completion of the previous one, mirroring the temporal dependency structure of RNNs. As a result, although Coconut operates within a Transformer architecture, its latent reasoning mechanism inherits the same computational bottlenecks associated with recurrent models, limiting training efficiency and scalability.

\paragraph{Computational Bottleneck.}
Because latent thoughts are generated one at a time, Coconut incurs a linear increase in computation with respect to the number of latent steps. This significantly increases training time, particularly when scaling to longer reasoning chains or larger models. Although key--value caching can mitigate redundant computation, the forward passes remain sequential and prevent efficient utilization of modern parallel hardware.

\paragraph{Implicit and Uncontrolled Latent Structure.}
Coconut relies on the emergent properties of hidden-state superposition to encode multiple potential reasoning paths within a single latent representation. While this can give rise to search-like behavior, the structure and content of latent reasoning are not explicitly controlled. The model must implicitly learn how to allocate representational capacity across possible reasoning branches, which can lead to instability and makes optimization more challenging.

\paragraph{Limited Scalability.}
Taken together, the sequential nature of latent generation, the lack of parallelism, and the dependence on curriculum-based supervision limit the scalability of Coconut. These constraints motivate the need for a reformulation that preserves continuous latent reasoning while enabling efficient parallel computation and stronger inductive structure.

\subsection{Overview of Codebook Continuous Thought}
\label{sec:overview_cct}

We introduce \emph{Codebook Continuous Thought} (CCT), a reformulation of continuous latent reasoning that replaces sequential latent unrolling with parallel latent synthesis. The core idea is to decouple the \emph{number of latent reasoning vectors} from the number of forward passes through the transformer by generating multiple latent representations simultaneously within a single reasoning iteration.

Instead of recursively feeding a single hidden state back into the model, CCT maintains a fixed-size set of learnable latent query vectors (a \emph{codebook}) that extract relevant information from the model’s recent hidden states via cross-attention. These latent vectors serve as continuous reasoning states and are generated in parallel. Within each iteration, the latents are allowed to interact bidirectionally with one another, forming a latent deliberation workspace, while interactions between text tokens and latents remain causally constrained.

At a high level, a single reasoning iteration in CCT follows the pattern
\[
\text{Text (causal)} \;\rightarrow\; \text{Latents (parallel, bidirectional)}
\]
\[
\text{Latents (parallel, bidirectional)} \;\rightarrow\; \text{Text (causal)}
\]
thereby preserving the autoregressive semantics of language modeling while enabling efficient and expressive latent-space reasoning.


\subsection{Parallel Latent Reasoning Architecture}
\label{sec:parallel_latent_reasoning}

We now present the full Codebook Continuous Thought (CCT) architecture, describing how latent reasoning states are initialized, synthesized, and refined within a single reasoning iteration. The design is guided by two principles: (i) latent reasoning should be generated efficiently and in parallel, and (ii) the resulting computation should preserve the causal semantics of autoregressive language modeling while providing sufficient flexibility for deliberative reasoning.

\paragraph{Latent Codebook and Initialization.}
CCT maintains a fixed-size set of $L$ learnable latent vectors, referred to as a \emph{latent codebook},
\[
\mathcal{Q} = \{ q_1, \dots, q_L \}, \quad q_i \in \mathbb{R}^d.
\]
These vectors act as persistent query slots that define the capacity of the latent reasoning workspace. At initialization, the codebook vectors are learned parameters and do not depend on the input sequence.

To provide a strong inductive bias and anchor latent reasoning to the task specification, we initialize the output projection layer of the cross attention block to zeros and use a residual connection from the hidden representation of the question. Let $h_q \in \mathbb{R}^d$ denote the hidden state corresponding to the final token of the question. After latent synthesis (described below), each latent vector receives a residual addition of $h_q$, ensuring that all latent reasoning states are explicitly conditioned on the question context. This initialization serves as a prior over latent reasoning and stabilizing training 

\paragraph{Parallel Latent Synthesis via Cross-Attention.}
Given the current textual context, let $h_t \in \mathbb{R}^d$ denote the hidden state at the final position, and define
\[
H_{t-k:t} = [h_{t-k}, \dots, h_t] \in \mathbb{R}^{(k+1) \times d}
\]
as the hidden states from the most recent $k+1$ positions. Latent reasoning states are synthesized in parallel using a cross-attention operation, where the latent codebook vectors act as queries and the recent hidden states serve as keys and values:
\[
Z = \mathrm{CrossAttn}\bigl(Q = \mathcal{Q},\; K = H_{t-k:t},\; V = H_{t-k:t}\bigr).
\]
This operation produces $Z \in \mathbb{R}^{L \times d}$, a set of $L$ continuous latent representations generated simultaneously in a single forward pass. In contrast to Coconut’s sequential latent unrolling, the number of latent reasoning states is no longer tied to the number of Transformer invocations, enabling efficient parallel computation.

% Restricting cross-attention to $H_{t-k:t}$ is essential for causal consistency. By conditioning latent synthesis only on a bounded window of past hidden states, the model enforces a Markovian update structure. As a result, the latent reasoning process remains compatible with standard autoregressive training objectives.

\paragraph{Latent Deliberation via Bidirectional Interaction.}
After parallel synthesis, the latent reasoning vectors are refined through a latent-only self-attention layer with a bidirectional attention mask. Before applying self-attention, we incorporate a residual connection from the hidden representation of the question to each latent vector. Let $h_q \in \mathbb{R}^d$ denote the hidden state corresponding to the final token of the question, and let $Z \in \mathbb{R}^{L \times d}$ denote the synthesized latent vectors. We form
\[
\tilde{Z} = Z + \mathbf{1} h_q^\top,
\]
where $\mathbf{1} \in \mathbb{R}^{L}$ is a vector of ones, ensuring that each latent receives the same residual contribution from $h_q$.

This residual grounding serves two purposes. First, it anchors all latent reasoning vectors to a shared semantic representation of the task specification, ensuring that deliberation remains explicitly conditioned on the question throughout the reasoning process. Second, it bounds the latent update to depend only on the current textual state, rather than the full history of latent computations, thereby reinforcing a Markovian structure in latent reasoning. Each iteration updates the latent state as a function of the current context and a fixed task representation, rather than recursively accumulating unstructured latent history.

The grounded latent representations $\tilde{Z}$ are then processed using bidirectional self-attention:
\[
Z' = \mathrm{SelfAttn}(\tilde{Z}),
\]
which allows latent vectors to exchange information freely within the latent workspace. This bidirectional interaction enables coordination and refinement among multiple candidate reasoning states, analogous to parallel hypothesis evaluation. Crucially, bidirectional attention is confined exclusively to the latent space within a single reasoning iteration: latent vectors do not attend to future text tokens, and text tokens do not attend to future latent vectors.

Together, these components define a unified latent reasoning module that replaces sequential continuous thought with an efficient, parallel, and causally consistent alternative, while preserving the expressive advantages of reasoning in continuous latent space.